\documentclass[11pt,a4paper]{article}
\usepackage[T1]{fontenc}
\usepackage{lmodern,microtype}
\usepackage[hmargin=25mm,vmargin=23mm]{geometry}
\usepackage{amsmath,amssymb,amsthm,booktabs}
\usepackage[hidelinks]{hyperref}
\hypersetup{
  pdftitle={Entropic Semantic Diffusion v2.0.0: Free wave amplitudes and bounded reconstruction gains},
  pdfauthor={Lorenzo Moriondo}}
\newcommand{\R}{\mathbb{R}}
\newcommand{\U}{U_q}
\DeclareMathOperator{\li}{li}
\DeclareMathOperator{\Ei}{Ei}
\DeclareMathOperator{\diag}{diag}
\DeclareMathOperator{\atanh}{atanh}
\newtheorem{proposition}{Proposition}
\title{\textbf{Entropic Semantic Diffusion v2.0.0}\\
  Free wave amplitudes and bounded reconstruction gains}
\author{Lorenzo Moriondo\\[-1pt]
  \small tuned.org.uk \quad\textbar\quad
  \href{https://orcid.org/0000-0002-8804-2963}{ORCID 0000-0002-8804-2963}}
\date{November 2026}

\begin{document}
\maketitle
\begin{abstract}
Projecting a wave state onto the unit sphere makes two signals in the same
direction indistinguishable even when their amplitudes differ. Version 2.0.0
of the Entropic Semantic Diffusion Model (ESDM) retains the raw state and
limits the gain applied to each retained graph mode at reconstruction. The
clock still advances by changes in the active modes' variational divergence.
For the readout, we compare the inclusive prime count with the logarithmic
integral at fixed integer arguments. Their ratio sets the initial gain;
a relative error envelope sets the interval in which that gain can be learned.
Schoenfeld's theorem motivates the interval under the Riemann hypothesis,
but enforcing it requires only a bounded function. We prove a relative norm
bound for the deterministic readout and give a scaling example with no
absolute amplitude bound. On 512 held-out noisy path signals, fitted bounded
gains reduced per-feature mean squared error from $0.022515$ (identity) to
$0.022373$; unrestricted gains reached $0.022320$. Shuffling the initial
prime residuals gave essentially the same error as keeping their order.
After 32 end-to-end updates, a linear probe recovered the inputs less
accurately from the encoder's posterior mean. These experiments establish
neither a semantic advantage for primes nor improved latent encoding.
\end{abstract}

\section{What the readout should preserve}
Consider two feature signals $u$ and $10u$, where $u$ is a unit eigenvector
of a graph Laplacian. Both have the same Rayleigh quotient. A unit-norm
projection also gives them the same state, so no decoder observing only that
state can recover which amplitude was present. A diagonal gain $w$, in
contrast, produces $wu$ and $10wu$. This difference matters even for
unit-normalised observations: a preprocessing convention need not constrain a
latent state after it has received a restoring force and external forcing from
another model component.

ESDM uses the same embedding in two ways. Its entries are coordinates of an
observed item; they are also a signal on a graph whose vertices are the
feature dimensions. ArrowSpace supplies this second geometry from a corpus
\cite{arrowspace}. ESDM inherits its spectral loading and wave block from
the Vibrational Deduction Transformer \cite{vdt}. A variational encoder
conditions the wiring, while an internal clock determines the wave step and the
amount of signal admitted by the diffusion decoder. The historical v0.4.1
manuscript describes a related image model, ArrowSpace Latent Diffusion with
Spectral Chart Conditioning (ALD-SC), that denoises a spatial latent
\cite{archive}. Here the object being
reconstructed is the feature signal of the entropic-clock ESDM. A bound on
this specific map cannot be transferred to an image decoder or to the entire
diffusion trajectory.

\section{From an observation to a wave state}
\paragraph{Geometry and encoding.}
Let an observation $x\in\R^{1\times F}$ be a row vector on a feature graph
with real symmetric Laplacian $L_F$. Retain $1\le q\le F$ orthonormal
eigenvectors in ascending eigenvalue order, as columns of
$\U=[u_1,\ldots,u_q]\in\R^{F\times q}$. Write $P=\U\U^\top$.
The frozen graph supplies the reference frame; the model has a separate
data-dependent wiring. Two encoder heads give a mean $\mu(x)$ and
log-variance $\ell(x)$ for a diagonal Gaussian:
\begin{equation}
 z=\mu(x)+\exp\!\bigl(\ell(x)/2\bigr)\odot\xi,\qquad
 \xi\sim\mathcal N(0,I_q).
\end{equation}
Numerical bounds on $\ell$ control scale parameters, not the realised
sample $z$. A spectral loading network predicts positive mode loadings
$\omega$ and a mixing matrix $S$. From
$W=\U\diag(\omega)S$ and the base edges it constructs an effective Laplacian
$L(z)$. Neither $W$ nor $\omega$ is the reconstruction gain defined below.

\paragraph{Two sectors and their clock.}
For each item, loadings strictly above the median across its modes form the
bright (active) set. The remaining modes, the latent wiring and the frozen
geometric prior form the dark sector. Let $K_n$ be the batch mean of the
sum of the active modes' Kullback--Leibler (KL) divergences at update $n$.
The spectral loading KL controls coupling to the prior. The split is
inspired by relational time in a cold-atom system \cite{barontini}; $K_n$
is a model statistic, not a measurement of thermodynamic entropy. Its total
variation defines the clock and a cosine signal schedule:
\begin{equation}
 \Delta\tau_n=|K_n-K_{n-1}|,\quad
 \tau_n=\sum_{j=1}^n\Delta\tau_j,\quad
 \bar\alpha(\tau_n)=\cos^2\!\left(
 \frac{\pi}{2}\min\{1,\tau_n/(3S_0)\}\right).
 \label{eq:clock}
\end{equation}
Here $S_0>0$ is the first-batch entropy budget. The clock is nondecreasing;
it does not advance between equal readings. Successive small increments can
trigger a stopping rule. The entropic pump $\Lambda_n$ estimates the inverse
magnitude of the finite-difference slope of $K$ against the log mean loading.
The implementation regularises the denominator, caps the pump at $0.1F$
before gating it by $\bar\alpha$, and limits the wave step $h$ by both
$\Delta\tau_n$ and a spectral stability condition.

\paragraph{The hybrid operator and recurrence.}
At a wave substep, let $L$ denote the symmetric effective Laplacian used by
the block (the base graph in the fixed-wiring case). First project the
current row $Q_t$ onto the reference chart and take one heat step. Evaluate
the Rayleigh gradient there and project the force onto the same chart:
\begin{align}
 f&=Q_t P,\qquad g=f-\tfrac12\sigma^2h\,fL,\qquad
 R_L(g)=\frac{gLg^\top}{gg^\top}, \label{eq:heat}\\
 \mathcal H_h(Q_t)&=\eta\,\frac{2(gL-R_L(g)g)}{gg^\top}\,P .
 \label{eq:hybrid}
\end{align}
The expression is row-wise for nonzero $g$; the implementation floors the
squared norm in the denominator. This is a gradient \emph{evaluated at}
the filtered signal, not a chain-rule derivative through the entire heat
step. With damping $\gamma$ and learned attention forcing $B_t$, the update
is
\begin{equation}
 Q_{t+1}=2Q_t-Q_{t-1}-\gamma(Q_t-Q_{t-1})
          -h^2\mathcal H_h(Q_t)+h^2\Lambda_n B_t .
 \label{eq:wave}
\end{equation}
A forward call starts from rest, $Q_{-1}=Q_0=x$: the wave receives the
observation directly, whereas $z$ conditions its wiring. At zero elapsed time the
raw state therefore remains $x$. No step projects $Q_{t+1}$ to unit length.
The clock and damping regulate the dynamics, but neither provides a uniform
amplitude bound. In particular, linear scaling of the \emph{readout} below
does not assume that the nonlinear recurrence scales linearly.

\section{How prime counts specify gains}
Use $s$ for an arithmetic argument, to distinguish it from a feature vector.
The right-continuous prime count is $\pi(s)=\#\{p\le s:p\text{ prime}\}$:
a prime at the evaluation point is counted in full, and no higher prime power
receives a contribution. For $s>1$ define the real, unoffset logarithmic
integral by \cite[Eq.~6.2.8]{dlmf}
\begin{equation}
 \li(s)=\operatorname{PV}\!\int_0^s\frac{dt}{\log t}=\Ei(\log s).
\end{equation}
The principal value is taken at $t=1$ and logarithms are natural.%
\footnote{The offset convention $\operatorname{Li}(s)=\int_2^s dt/\log t
=\li(s)-\li(2)$ differs by $\li(2)\simeq1.04516$. Substituting it into
the residual or the ratio changes the initialisation.}
Under the classical Riemann hypothesis (RH), which places all nontrivial
zeros of the zeta function on the line $\Re s=1/2$,
Schoenfeld's Corollary~1 \cite[Eq.~(6.18)]{schoenfeld} gives
\begin{equation}
 |\pi(s)-\li(s)| < B(s),\qquad
 B(s)=\frac{\sqrt{s}\log s}{8\pi},\qquad s\ge2657 .
 \label{eq:rh}
\end{equation}
The $\pi$ in $8\pi$ is the circle constant. Equation~\eqref{eq:rh} concerns
an error in counting primes. It makes no statement about a Laplacian,
an embedding or a wave amplitude. For the finite interval we use, the
same arithmetic inequality is also known without assuming full RH
\cite{johnston}.

Assign to column $u_k$ the deterministic grid point
\begin{equation}
 s_k=\operatorname{round}\!\left(4096\cdot
          16^{(k-1)/(q-1)}\right),\quad k=1,\ldots,q,\quad q>1;
 \qquad s_1=4096\quad(q=1).
\end{equation}
The endpoints $4096$ and $65536$ span four doublings above the theorem's
threshold. The assignment is by column index, not by the numerical value of
a graph frequency. It is a modelling choice, not a consequence of the
prime number theorem. Fix the half-width $d_k$ and compute the scaled
arithmetic residual $r_k$ at each point:
\begin{align}
 d_k&=\frac{B(s_k)}{\li(s_k)}, &
 r_k&=\frac{\pi(s_k)-\li(s_k)}{B(s_k)}, \label{eq:width}\\
 \theta_k^{(0)}&=\atanh(r_k), &
 w_k(\theta_k)&=1+d_k\tanh\theta_k . \label{eq:gain}
\end{align}
Verify $|r_k|<1$ before taking the inverse hyperbolic tangent; clipping a
failed check would conceal an incorrect count, convention or grid. The
initial gain then equals $\pi(s_k)/\li(s_k)$. The $q$ logits $\theta_k$
can be optimised independently. After initialisation there is no requirement
for a learned gain to equal the prime-count ratio.

For example, at $4096$ there are $564$ primes, $\li(4096)\simeq576.923$,
and the initial gain is $0.977601$ with half-width $0.036714$.
At $65536$ these values are $6542$, $6583.987$, $0.993623$ and
$0.017158$, respectively. An inclusive sieve and evaluation of $\Ei(\log s)$
over every integer from $4096$ through $65536$ give
$\max_s |(\pi(s)-\li(s))/B(s)|\simeq0.825800$ and
$\max_s B(s)/\li(s)\simeq0.036714$. These finite numerical checks concern
the initialiser; the parameterisation itself enforces the gain interval
independently of RH at every subsequent update.

\section{What the reconstruction bound implies}
Let $Q$ denote the final raw wave row, $a=Q\U$ its retained coordinates and
$Q_\perp=Q(I-P)$ its complement. The deterministic signal readout is
\begin{align}
 D_w(Q)&=Q+\bigl((Q\U)\odot(w-\mathbf1)\bigr)\U^\top
       =Q_\perp+(a\odot w)\U^\top. \label{eq:readout}
\end{align}
The diffusion decoder forms
\begin{equation}
 \hat x=\sqrt{\bar\alpha(\tau)}\,D_w(Q)
          +\sqrt{1-\bar\alpha(\tau)}\,\epsilon .
 \label{eq:recon}
\end{equation}
Forward reconstruction sets $\epsilon=0$; in a stochastic decode the
additive Gaussian noise is outside $D_w$. In particular, $\bar\alpha(0)=1$ gives
the learned readout, not necessarily an identity reconstruction. All
unretained directions pass through the signal readout with gain one; the
subsequent clock factor multiplies both retained and unretained signal.
The reverse sampler also uses $D_w$ at its signal readouts. Its initial
wave and injected noise have Gaussian coordinates in the frozen chart;
restricting their directions does not fix their magnitudes.

\begin{proposition}[Relative distortion of the signal]
For an orthonormal chart $\U^\top\U=I_q$ and finite real logits, set
$d_*=\max_k d_k$. Then $|w_k-1|<d_k$ and
\begin{equation}
 \|D_w(Q)-Q\|_2\le d_*\|Q\U\|_2\le d_*\|Q\|_2 .
 \label{eq:bound}
\end{equation}
Consequently $(1-d_*)\|Q\|_2\le\|D_w(Q)\|_2\le(1+d_*)\|Q\|_2$.
\end{proposition}
\begin{proof}
The first assertion follows from $|\tanh\theta|<1$. For any chart row $b$,
$\|b\U^\top\|_2^2=b\U^\top\U b^\top=\|b\|_2^2$. Therefore
\[
 \|D_w(Q)-Q\|_2^2=\sum_{k=1}^q a_k^2(w_k-1)^2
       \le d_*^2\|a\|_2^2.
\]
Orthogonality gives $\|Q\|_2^2=\|Q_\perp\|_2^2+\|a\|_2^2$.
Taking square roots gives \eqref{eq:bound}; the triangle and reverse triangle
inequalities give the norm bounds. Floating-point saturation of $\tanh$ can
reach an endpoint, so numerical checks allow the closed interval to rounding
precision.
\end{proof}

For a fixed readout $D_w(cu_k^\top)=cw_k u_k^\top$ for any real $c$.
Since $d_*<1$ on the chosen grid, $w_k$ is nonzero and the output norm grows
without limit with $|c|$. A row orthogonal to every $u_k$ is reproduced
exactly. Thus this readout preserves amplitude information even though it
may attenuate or amplify a retained component. It does not itself compress
an $F$-vector to $q$ numbers: the complement would be lost in a chart-only
encoder, but is transmitted in \eqref{eq:readout}.

Changing the sign of an eigenvector changes neither its rank-one projector
nor the readout. Within a repeated-eigenvalue subspace the choice of
orthonormal frame is arbitrary. Rotating that frame generally changes the
readout if gains differ; even the retained subspace may depend on the frame
if the truncation splits a repeated eigenvalue. Equation~\eqref{eq:bound}
holds for each fixed frame. A checkpoint must keep its basis and gain order
together to reproduce the same operator.

\paragraph{Training objective.}
In loss-minimisation form the regularised negative evidence lower bound
(ELBO) used by ESDM is
\begin{equation}
 \mathcal J=\operatorname{MSE}(x,\hat x)+K_z+K_S+K_n+
      \beta(K_S+K_n+S_{\rm dark}-S_0)^2 .
\end{equation}
Here MSE is mean squared error. $K_z$ regularises the diagonal Gaussian
against a standard normal; $K_S$
regularises the spectral loading posterior; $K_n$ is the bright KL that
also advances the clock. The implementation uses $\beta\simeq1/S_0$.
After the first batch $S_{\rm dark}=\max(0,S_0-K_S-K_n)$, so whenever
$K_S+K_n\le S_0$ the last term vanishes by construction. It cannot by
itself demonstrate physical conservation. The gains enter the objective
through the reconstruction, without a prime-count penalty in the ELBO.

\section{Synthetic encoding and decoding}
\paragraph{Experimental design.}
We draw 512 training and 512 independent test signals on an unweighted path
with $F=32$ vertices (CPU, seed 3407). In its orthonormal discrete cosine
basis the clean coefficients are independent zero-mean Gaussians with
variance $v_j=1/(1+(j/4)^2)$, $j=0,\ldots,31$. Thus smooth modes are more
energetic by design; the signal distribution makes no reference to primes.
The decoding experiment adds independent Gaussian observation noise with
standard deviation $0.15$ to each feature. Errors are per-feature squared
differences, averaged within each held-out case and then over cases.

\paragraph{Encoding.}
At ranks $q=2,4,8,16,32$, we encode \emph{clean} held-out signals by
$x\U$ and synthesise with $\U^\top$. We compare this chart-only projection
with two codes from the untrained ESDM Gaussian encoder at the same rank:
the posterior mean and one posterior draw per item. A separate affine
least-squares probe, fitted only on training codes and clean targets,
decodes each type of Gaussian code. The chart needs no fit. These
comparisons measure recoverability from the code alone; because $Q_0=x$,
a forward pass through the entire model is not a bottleneck test.

\paragraph{Decoding.}
At $q=8$ the input $Q$ is the noisy feature row, the target is the clean row,
and the clock and additive decoder noise are omitted to isolate $D_w$.
We compare five arms: identity held fixed; bounded logits initialised at
zero (``smooth''); bounded logits initialised by the observed residuals
(``prime''); the same normalised residuals permuted among columns, retaining
the original $d_k$ (``shuffled''); and unrestricted diagonal gains
initialised at one (``free''). The last four arms receive 100 full-batch Adam
updates at learning rate $0.05$ on the same training pairs, with no early
stopping. All five keep the complement at unit gain. We compare errors
on identical test cases; 95\% percentile intervals for paired differences
use 2,000 case bootstrap draws, conditional on this split and fit. They
do not measure uncertainty across seeds or training runs.

\paragraph{What the controls predict.}
For a clean scalar coefficient $a$ of variance $v_j$ and independent noise
$e$ of variance $\sigma_e^2=0.15^2$, the expected error of a linear gain is
\[
 \mathbb E[(w(a+e)-a)^2]=(w-1)^2v_j+w^2\sigma_e^2.
\]
The minimiser is $v_j/(v_j+\sigma_e^2)$. At $j=0$ this is $0.9780$, close
to the first prime-initialised gain ($0.9776$). For modes $j=3$ through
$7$ the optimum lies below the engineered lower bound: at $j=7$ the
optimal gain is $0.91625$ but the interval starts at $0.98284$.
An unrestricted gain can therefore remove noise that the
bounded gain must transmit. All arms also transmit noise on 24 unretained
directions, contributing $24\sigma_e^2/32=0.016875$ in expectation to
per-feature MSE. This term accounts for three quarters of the identity
arm's expected error and cancels in the paired contrasts.

\section{Results}
\paragraph{What survives encoding?}
Table~\ref{tab:encoding} shows that the fixed path chart recovered more of
the clean signal than either code from the untrained encoder at low ranks.
This is expected: the chart is the generator's variance-ordered basis.
At rank eight, replacing the mean by one posterior sample increased probe
MSE by $0.10294$ (paired 95\% interval $[0.09777,0.10856]$). At full rank
both the orthonormal chart and the probed affine mean map were numerically
invertible on these signals. The remaining error of the sampled code does
not measure the performance of a trained variational encoder.

\begin{table}[htbp]
\centering
\caption{Encoding: held-out mean squared error per original feature for 512
clean signals. The chart uses transpose synthesis; each posterior code uses
its own affine probe fitted on 512 training cases. Lower is better.}
\label{tab:encoding}
\begin{tabular}{@{}rrrr@{}}
\toprule
Rank $q$ & Path chart & Posterior mean & One posterior sample\\
\midrule
2  & 0.13686 & 0.15640 & 0.19246\\
4  & 0.09269 & 0.13528 & 0.19358\\
8  & 0.04536 & 0.08821 & 0.19115\\
16 & 0.01563 & 0.04326 & 0.19046\\
32 & $1.47\times10^{-30}$ & $1.03\times10^{-13}$ & 0.18102\\
\bottomrule
\end{tabular}
\end{table}

\paragraph{Does end-to-end training improve the code?}
We trained one rank-eight ESDM on the 512 clean training signals, using
the full ELBO loss for two epochs (seed 3407, minibatches of 32, Adam
learning rate $0.0003$, weight decay $0.00001$). All 32 updates were finite;
the largest absolute change to a readout logit was $0.010664$. We fitted
separate affine probes to the training-set posterior means before and after
training. On the same 512 held-out cases, probe MSE rose from $0.08821115$
to $0.09270693$, a paired increase of $0.00449578$ (95\% case-bootstrap
interval $[0.00398528,0.00502858]$). The posterior mean became less useful
to this probe in this run. Because the wave path also receives the original
features, this result does not measure the change in full-model reconstruction.

\paragraph{Does the order of prime residuals help?}
At initialisation the prime arm improved on identity, but shuffling the
residuals preserved that improvement (Table~\ref{tab:decoding}). After
fitting, the MSE difference for prime minus shuffled was
$1.31\times10^{-7}$ (95\% interval $[-0.79,3.28]\times10^{-7}$).
The prime arm finished slightly below the smooth start
($-7.94\times10^{-7}$, interval $[-15.58,-0.41]\times10^{-7}$).
All three bounded arms have the same feasible set; differences after a
fixed optimisation budget need not persist at convergence. The unrestricted
diagonal achieved the lowest error: prime minus free was
$5.34\times10^{-5}$, interval $[2.35,8.45]\times10^{-5}$.

\begin{table}[htbp]
\centering
\caption{Decoding: held-out per-feature MSE for identical noisy inputs and
clean targets ($q=8$, 512 test cases). Min/max are over the eight final
gains, not reconstructed coordinates. Identity remains fixed.}
\label{tab:decoding}
\begin{tabular}{@{}lrrrr@{}}
\toprule
Readout & Initial MSE & Final MSE & Min.\ gain & Max.\ gain\\
\midrule
Identity           & 0.022514806 & 0.022514806 & 1.00000 & 1.00000\\
Smooth start       & 0.022514806 & 0.022374183 & 0.97089 & 0.98329\\
Prime start        & 0.022420285 & 0.022373389 & 0.97055 & 0.98312\\
Shuffled residuals & 0.022414795 & 0.022373259 & 0.97040 & 0.98306\\
Free diagonal      & 0.022514806 & 0.022319945 & 0.94371 & 0.97983\\
\bottomrule
\end{tabular}
\end{table}

\paragraph{Where does the model apply the bound?}
In the actual \texttt{ESDM.readout}, input coefficients $1$, $10$ and $1000$
on the first mode gave $0.977601$, $9.77601$ and $977.601$. We also ran two
consecutive clean batches of eight through the untrained model. For the
first, $\tau=0$ and $\bar\alpha=1$; the raw wave norms ranged from $1.815$
to $2.715$. For the second, $\tau=3.2323$ and $\bar\alpha=0.99598$;
the raw norms ranged from $2.136$ to $3.367$. Reapplying
$\sqrt{\bar\alpha}D_w$ to the stored pre-readout state reproduced the model
output with maximum absolute discrepancy zero in both passes. This checks
the placement of the bound, not the quality of a trained model.

\section{Interpretation and reproducibility}
The guaranteed property is a small, mode-dependent change to each signal
relative to its input. The empirical trade-off on this path was a small
improvement over identity at the cost of higher error than the free diagonal.
The prime and shuffled starts do not establish a benefit specific to prime
ordering. One graph, one split and one initialisation cannot establish a
semantic or generative advantage, or validate long-run wave stability. The
prime-grid calculation is finite and supplies no evidence for RH.

To reproduce Tables~\ref{tab:encoding} and~\ref{tab:decoding}, run the
following command from the repository root \cite{archive}:
\begin{verbatim}
uv run --no-sync python notebooks/run.py
\end{verbatim}
\noindent The runner executes \path{notebooks/encoding-quality.ipynb} and
\path{notebooks/decoding.ipynb}. The full-precision results, split hashes,
software versions and forward checks are in
\path{docs/paper/v2-results/experiment.json}.

\clearpage
\begin{thebibliography}{9}
\bibitem{arrowspace}
L.~Moriondo. ArrowSpace: introducing Spectral Indexing for vector search.
\emph{J. Open Source Softw.} 10, 9002 (2025).
\href{https://doi.org/10.21105/joss.09002}{doi:10.21105/joss.09002}.
\bibitem{vdt}
L.~Moriondo. Vibrational Deduction Transformer (VDT) Leveraging Spectral
methods, Theory of Sound and Associative Memory. Zenodo (2026).
\href{https://doi.org/10.5281/zenodo.20816835}{doi:10.5281/zenodo.20816835}.
\bibitem{archive}
L.~Moriondo. \emph{Entropic Semantic Diffusion: ALD-SC---A Spectral Latent
Diffusion Model}. July 2026 manuscript preserved as
\texttt{docs/paper/esdm.tex}, v0.4.1, in the
\href{https://github.com/tuned-org-uk/entropic-semantic-diffusion}{ESDM repository}.
\bibitem{barontini}
G.~Barontini. Testing the problem of time with cold atoms.
\emph{Phys. Rev. Research} 8, L022047 (2026).
\href{https://doi.org/10.1103/1h9j-df4k}{doi:10.1103/1h9j-df4k}.
\bibitem{dlmf}
NIST Digital Library of Mathematical Functions.
\href{https://dlmf.nist.gov/6.2.E8}{Eq.~6.2.8, Logarithmic integral}.
\bibitem{schoenfeld}
L.~Schoenfeld. Sharper bounds for the Chebyshev functions $\theta(x)$ and
$\psi(x)$. II. \emph{Math. Comp.} 30, 337--360 (1976), Cor.~1, p.~339.
\href{https://doi.org/10.1090/S0025-5718-1976-0457374-X}{doi:10.1090/S0025-5718-1976-0457374-X}.
\bibitem{johnston}
D.~R.~Johnston. Improving bounds on prime counting functions by partial
verification of the Riemann hypothesis. \emph{Ramanujan J.} 59, 1307--1321
(2022). \href{https://doi.org/10.1007/s11139-022-00616-x}{doi:10.1007/s11139-022-00616-x}.
\end{thebibliography}
\end{document}
